\section{Bernhard Riemann} % lang-pl
\label{sekcja_riemann}

Hipotezę kąta rozwartego wszyscy dotąd odrzucali, bo przeczyła nieskończoności prostej. % lang-pl
W tej sekcji zobaczymy, jak uczeń Gaussa odróżni nieograniczoność od nieskończoności i jak jeden wykład otworzy drugi okres w dziejach geometrii nieeuklidesowej. % lang-pl

Bernhard Riemann (1826--1866) urodzi się w Breselenz w Królestwie Hanoweru, będzie studiował w Getyndze u Gaussa, który zostanie jego promotorem i zachęci go do matematyki zamiast teologii. % lang-pl
Zachoruje na gruźlicę i umrze we Włoszech w wieku trzydziestu dziewięciu lat. % lang-pl
\index[persons]{Riemann, Bernhard}%

Wykład habilitacyjny \emph{Über die Hypothesen, welche der Geometrie zu Grunde liegen} wygłosi 10 czerwca 1854 roku w Getyndze; wydrukuje go dopiero w 1868 roku, już po śmierci Riemanna, Richard Dedekind. % lang-pl
Riemann omówi w nim pojęcia nieograniczoności i nieskończoności; gdy różnica między nimi stanie się jasna, da się zrealizować równie niesprzeczną geometrię hipotezy kąta rozwartego, jeśli zmodyfikuje się postulaty Euklidesa \cite[s. 287]{eves1_1972}: % lang-pl
\index[persons]{Dedekind, Richard}%
\begin{itemize}
\item dwa różne punkty wyznaczają co najmniej jedną prostą; % lang-pl
\item prosta jest nieograniczona; % lang-pl
\item każde dwie proste na płaszczyźnie przecinają się. % lang-pl
\end{itemize}
Eves zauważa, że pierwsze dwa z tych zdań to po prostu postulaty Euklidesa wzięte dosłownie: Euklides miał na myśli, że dwa punkty wyznaczają \emph{dokładnie} jedną prostą i że jest ona \emph{nieskończona}, ale tego nie napisał \cite[s. 288]{eves1_1972}. % lang-pl

Wykład Riemanna otworzy drugi okres rozwoju geometrii nieeuklidesowej: zamiast metod elementarnej geometrii syntetycznej stosować się będzie metody geometrii różniczkowej. % lang-pl
Riemann stworzy w ten sposób całą klasę nietradycyjnych geometrii, zwanych \emph{riemannowskimi}, a uogólnienie pojęcia przestrzeni, do którego prowadzi jego wykład, zastosuje się później w teorii względności \cite[s. 288]{eves1_1972}. % lang-pl
Eugenio Beltrami powoła się na ten wykład w swoich pracach o modelach (zobacz sekcję \ref{sekcja_beltrami_cayley_klein}), które ukażą się w tym samym 1868 roku, w którym Dedekind wydrukuje wykład. % lang-pl
\index{geometria!riemannowska} % lang-pl

Eves wspomina jeszcze inną nietradycyjną geometrię: Max Dehn zbuduje taką, w której zawieszony jest aksjomat Archimedesa (geometria nieArchimedesowa) \cite[s. 288]{eves1_1972}. % lang-pl
\index[persons]{Dehn, Max}%

\todofoot{Geometria riemannowska (metryka, krzywizna) wykracza poza Evesa -- Eves \cite[s. 315]{eves1_1972} sam jej nie rozwija; potrzebne inne źródło, jeśli chcemy ją opisać.} % lang-pl
% Źródła: Eves s. 289--290 (zad. 10), 314--315; Wikipedia: Girard's theorem (Spherical trigonometry), Lexell's theorem.
% https://en.wikipedia.org/wiki/Lexell's_theorem
% https://en.wikipedia.org/wiki/Spherical_trigonometry

